\documentclass[12pt,letterpaper]{article}
\usepackage[utf8]{inputenc}
\usepackage[spanish]{babel}
\usepackage{amsmath}
\usepackage{amsfonts}
\usepackage{amssymb}
\usepackage[margin=2.5cm]{geometry}

\begin{document}

\begin{center}
    \Large \textbf{Evaluación de Examen 2: Unidad II} \\
    \large Física para Ciencias de la Salud (005-1713) \\
    \vspace{0.5cm}
    \textbf{Estudiante:} Thomairys Sosa \quad \textbf{C.I.:} 33.685.285
    \vspace{0.3cm} \\
    \large \textbf{Calificación Final: 10/10 puntos}
\end{center}

\vspace{0.5cm}

\section*{Análisis de Resultados}

\subsection*{1. Cinemática (Aceleración media)}
\textbf{Procedimiento y resultado:} Extrae de manera estructurada los datos iniciales y plantea la ecuación de aceleración media ($a_m = \frac{V_f - V_i}{T}$). Efectúa la sustitución correcta restando las velocidades y dividiendo entre el tiempo ($0,6 \text{ m/s} / 0,15 \text{ s}$). El cálculo aritmético es exacto y obtiene $4 \text{ m/s}^2$ con un correcto manejo dimensional. \\
\textbf{Veredicto:} \textbf{Correcto} (2/2 pts).

\subsection*{2. Dinámica (Leyes de Newton)}
\textbf{Procedimiento y resultado:} Plantea la fórmula original de la Segunda Ley de Newton ($F = m \cdot a$) y realiza el despeje impecable para hallar la aceleración ($a = \frac{f}{m}$). Sustituye los valores dividiendo $150 \text{ N}$ entre $25 \text{ kg}$, logrando el resultado perfecto de $6 \text{ m/s}^2$. \\
\textbf{Veredicto:} \textbf{Correcto} (2/2 pts).

\subsection*{3. Trabajo Mecánico}
\textbf{Procedimiento y resultado:} Identifica la fórmula general del trabajo. Efectúa correctamente la multiplicación ($400 \text{ N} \cdot 0,8 \text{ m}$) e indica en un excelente paso intermedio que esto equivale a $320 \text{ N} \cdot \text{m}$, para concluir de manera perfecta que el resultado es $320 \text{ Joules}$. \\
\textbf{Veredicto:} \textbf{Correcto} (2/2 pts).

\subsection*{4. Energía Potencial}
\textbf{Procedimiento y resultado:} Extrae rigurosamente los datos del problema (masa, altura y constante gravitacional). Plantea el producto de las tres variables según la ecuación de la energía potencial gravitatoria ($E_p = m \cdot g \cdot h$) e integra las unidades en su sustitución. Obtiene el resultado verificado y exacto de $17,64 \text{ Joules}$. \\
\textbf{Veredicto:} \textbf{Correcto} (2/2 pts).

\subsection*{5. Potencia Mecánica}
\textbf{Procedimiento y resultado:} Extrae los datos y relaciona correctamente el trabajo efectuado y el tiempo mediante la respectiva fracción ($P = \frac{W}{T}$). El desarrollo de la división de $75 \text{ J} / 60 \text{ s}$ arroja la cifra correcta y verificada de $1,25 \text{ Watts}$. \\
\textbf{Veredicto:} \textbf{Correcto} (2/2 pts).

\vspace{0.5cm}
\section*{Comentarios Generales y Sugerencias}
¡Felicidades por obtener una calificación perfecta! El examen demuestra un nivel de excelencia y madurez analítica notable.
\begin{itemize}
    \item El formato implementado en cada problema (separando de forma impecable la zona de \textit{Datos} y el desarrollo paso a paso de las ecuaciones) facilita enormemente la corrección y el seguimiento metodológico.
    \item El manejo de las magnitudes y el \textbf{análisis dimensional} es uno de los mejores evaluados, destacando la derivación paso a paso de unidades equivalentes (como los $\text{N} \cdot \text{m}$ a Joules en el problema 3).
    \item Se incentiva a la estudiante a mantener este estupendo estándar de dedicación, orden y limpieza en sus futuras evaluaciones y reportes.
\end{itemize}

\end{document}